\subsection{§III.1.9 (Prop.~4) --- $\inf A \le \sup A$ pour $A$ non vide}
\textbf{Énoncé (Bourbaki).} Soit $A \subset E$ une partie non vide d'un ensemble ordonné, admettant une borne inférieure $\inf A$ et une borne supérieure $\sup A$. Alors $\inf A \le \sup A$. L'hypothèse $A \neq \emptyset$ est essentielle : pour $A = \emptyset$ tout élément est à la fois minorant et majorant, et l'inégalité s'inverse.

\textbf{Formalisé.} Cible : $(i,s) \in G$, avec la convention de graphe $x \le y := (x,y)\in G$, $i = \inf A$, $s = \sup A$. \\
Module : \texttt{bourbaki/ordre/iii\_1\_relations\_ordre/bornes\_sup/ensembles\_inf\_sup\_prop4\_iii1.py}.

\textbf{Preuve (\Nstar).} On \texttt{assume} les quatre hypothèses puis, par \texttt{conjonction\_elim} sur les prédicats \texttt{borne\_inferieure}/\texttt{borne\_superieure}, on extrait que $i$ minore $A$ ($(\forall x)(x\in A \Rightarrow (i,x)\in G)$) et que $s$ majore $A$. Sous un témoin $a \in A$ (\texttt{assume}), \texttt{instancie} + \texttt{modus\_ponens} donnent $(i,a)\in G$ et $(a,s)\in G$ ; la transitivité instanciée en $(i,a,s)$ conjuguée par \texttt{conjonction\_intro} conclut $(i,s)\in G$. Le témoin est déchargé par \texttt{loi\_deduction} puis \texttt{existe\_elimination} : la conclusion ne contient pas $a$ libre, donc l'hypothèse honnête reste $(\exists z)(z\in A)$, pas le témoin.

\textbf{Statut.} CLOS sous hypothèses honnêtes $\{\,$\texttt{transitivite\_rel(G)}, $(\exists z)(z\in A)$, \texttt{borne\_inferieure(G,A,i,E)}, \texttt{borne\_superieure(G,A,s,E)}$\,\}$ ; \texttt{theorie\_ensembles()==22}. Les deux hypothèses d'existence des bornes reproduisent fidèlement le « lorsque les bornes existent » de Bourbaki.

\subsection{§III.1.10 (Prop.~10) --- Maximal $\Rightarrow$ plus grand dans un filtrant à droite}
\textbf{Énoncé (Bourbaki).} Soit $E$ un ensemble ordonné filtrant à droite (deux éléments quelconques admettent un majorant commun) et $a$ un élément maximal de $E$. Alors $a$ est le plus grand élément de $E$.

\textbf{Formalisé.} Cible : $\texttt{plus\_grand\_element}(G,E,a)$, c.-à-d. $a\in E \;\wedge\; (\forall x)(x\in E \Rightarrow (x,a)\in G)$. \\
Module : \texttt{bourbaki/ordre/iii\_1\_relations\_ordre/iii\_1\_8\_filtrants/ensembles\_prop10\_maximal\_filtrant.py}. Ce module comble le dossier-trou \texttt{iii\_1\_8\_filtrants}.

\textbf{Preuve (\Nstar).} On \texttt{assume} \texttt{est\_ordre}, \texttt{filtrant\_droite\_G} et \texttt{element\_maximal}, dont on extrait par \texttt{conjonction\_elim} : $a\in E$, le corps de maximalité, la réflexivité et la transitivité. Pour $x \in E$ (\texttt{assume}), le filtrant instancié en $(x,a)$ fournit par \texttt{modus\_ponens} un majorant commun $(\exists z)(z\in E \wedge (x,z)\in G \wedge (a,z)\in G)$. Sous le témoin $z$, la maximalité appliquée à $z$ (avec $(a,z)\in G$) force $z=a$ ; la réflexivité $(a,a)\in G$ est transportée par $a=z$ via Leibniz \texttt{N.s6} + \texttt{equivalence\_avant} en $(z,a)\in G$, puis la transitivité en $(x,z,a)$ donne $(x,a)\in G$. \texttt{existe\_elimination} décharge le témoin, \texttt{generalisation} (sur le liant interne fixé) ferme le corps « plus grand ». C'est ici que \texttt{est\_ordre} devient load-bearing (réflexivité + transitivité).

\textbf{Statut.} CLOS sous hypothèses honnêtes $\{\,$\texttt{est\_ordre(G,E)}, \texttt{filtrant\_droite\_G(G,E)}, \texttt{element\_maximal(G,E,a)}$\,\}$ ; \texttt{theorie\_ensembles()==22}. Une assertion interne vérifie \texttt{conclusion == cible}.

\subsection{§III.1.13 (Prop.~13) --- Intersection de deux intervalles fermés}
\textbf{Énoncé (Bourbaki).} L'intersection $\intEE{a}{b} \cap \intEE{c}{d}$ de deux intervalles fermés (même support $E$, même ordre $G$) est encore un intervalle fermé. La preuve « treillis » l'identifie à $\intEE{\sup(a,c)}{\inf(b,d)}$, ce qui suppose l'existence des bornes des paires.

\textbf{Formalisé.} Cible : $(\forall x)\bigl(x \in \intEE{a}{b}\cap\intEE{c}{d} \iff (x\in E \wedge (a,x)\in G \wedge (c,x)\in G \wedge (x,b)\in G \wedge (x,d)\in G)\bigr)$ (forme membership). \\
Module : \texttt{bourbaki/ordre/iii\_1\_relations\_ordre/ordre\_treillis/ensembles\_intervalles\_prop13.py}.

\textbf{Preuve (\Nstar).} L'axiome de membership de $\intEE{\cdot}{\cdot}$ vit dans la théorie dédiée \texttt{theorie\_intervalle\_ferme} et est déchargé par \texttt{N.axiome}, puis \texttt{instancie} en $(E,a,b,x)$ et $(E,c,d,x)$. Le lemme \texttt{\_instance\_intersection} (via \texttt{AXIOME\_INTER} de \texttt{theorie\_ensembles}) donne $x\in X\cap Y \iff (x\in X \wedge x\in Y)$. Sens $\Rightarrow$ : \texttt{assume} $x\in\intEE{a}{b}\cap\intEE{c}{d}$, puis \texttt{equivalence\_avant} + \texttt{modus\_ponens} + \texttt{conjonction\_elim} extraient $x\in E$, $a\le x$, $c\le x$, $x\le b$, $x\le d$, recombinés par \texttt{conjonction\_intro}. Sens $\Leftarrow$ : on reconstruit les deux memberships et on réinjecte par \texttt{equivalence\_arriere} dans les axiomes puis dans l'intersection. \texttt{conjonction\_intro} des deux sens donne l'équivalence, \texttt{generalisation} en $x$ ferme.

\textbf{Statut.} CLOS (0 hypothèse) : l'unique axiome de $\intEE{\cdot}{\cdot}$ est interne à \texttt{theorie\_intervalle\_ferme} (S8+A1) et l'intersection emploie l'\texttt{AXIOME\_INTER} déjà compté ; \texttt{theorie\_ensembles()==22}. Résidu honnête documenté : la forme treillis $\intEE{a}{b}\cap\intEE{c}{d}=\intEE{\sup(a,c)}{\inf(b,d)}$ n'est pas prouvée (elle exige l'existence de $\sup$/$\inf$ de paires) ; la forme membership ci-dessus est la caractérisation order-théorique complète et exacte.

\subsection{§III.1.5 --- Proposition 2 : connexion de Galois}
\textbf{Énoncé (Bourbaki, E~III.7--8).} Pour $u,v$ décroissantes avec $v\circ u\geq\mathrm{id}$
et $u\circ v\geq\mathrm{id}$ : $u\circ v\circ u = u$ (et $v\circ u\circ v = v$).

\textbf{Formalisé.} Première égalité $u\circ v\circ u=u$ : $\;(\forall x)(x\in E\Rightarrow w(x)=u(x))$,
$w$ représentant le composite. Module :
\texttt{.../iii\_1\_5\_applications\_croissantes/ensembles\_prop2\_galois.py}.

\textbf{Preuve (\Nstar).} $u(v(u(x)))\leq u(x)$ ($u$ décroissante appliquée à $v(u(x))\geq x$)
et $u(x)\leq u(v(u(x)))$ (Galois en $x'{:=}u(x)$) ; antisymétrie ; transport par
\texttt{composer\_egalites}.

\textbf{Statut.} CLOS sous $7$ hypothèses honnêtes. \emph{Fidélité} : la décroissance
de~$v$, non load-bearing pour cette égalité, a été retirée (cf.\ \texttt{DECISIONS.md}).

\subsection{§III.1.7 --- Le plus petit élément est l'unique élément minimal}
\textbf{Énoncé (Bourbaki).} Si $E$ admet un plus petit élément $a$, tout élément
minimal $m$ de $E$ coïncide avec $a$.

\textbf{Formalisé.} Cible : $m=a$, sous
$\{\mathrm{plus\_petit\_element}(G,E,a),\ \mathrm{element\_minimal}(G,E,m)\}$. \\
Module : \texttt{.../iii\_1\_relations\_ordre/iii\_1\_7\_plus\_grand\_plus\_petit/ensembles\_plus\_petit\_unique\_minimal.py}.

\textbf{Preuve (\Nstar).} $a$ minore $E$ instancié en $m$ donne $(a,m)\in G$ ; la
minimalité de $m$ instanciée en $a$ donne $a=m$ ; \texttt{symetrie}. L'antisymétrie
n'est pas load-bearing et n'est donc pas une hypothèse.

\textbf{Statut.} CLOS sous hypothèses honnêtes ; \texttt{theorie\_ensembles()==22}.

\subsection{§III.1.12 --- Proposition 12 : critère de borne supérieure (ordre total)}
\textbf{Énoncé (Bourbaki, E~III.14).} $E$ totalement ordonné, $X\subset E$, $b\in E$ :
$b=\sup X$ ssi (1°) $b$ majore $X$ et (2°) tout $c<b$ est dépassé par un $x\in X$ avec $c<x\leq b$.

\textbf{Formalisé.} L'équivalence $\mathrm{borne\_sup}(G,X,b,E)\Leftrightarrow(\text{1°}\wedge\text{2°})$,
les deux sens. Module :
\texttt{.../iii\_1\_12\_totalement\_ordonnes/ensembles\_prop12\_sup\_total.py}.

\textbf{Preuve (\Nstar).} Totalité $+$ antisymétrie (gabarit
\texttt{maximal\_est\_plus\_grand\_si\_total}) ; sens~$\Rightarrow$ par contraposée,
sens~$\Leftarrow$ par témoin. \emph{Fidélité} : le strict $c<x$ ($c\neq x$) est
\emph{obligatoire} --- sans lui l'équivalence est fausse (contre-exemple
$E=\{0,1,2\}$) et le noyau la refuse (cf.\ \texttt{DECISIONS.md}).

\textbf{Statut.} CLOS sous $3$ hypothèses honnêtes.
